\chapter{Conception de l'architecture multi-agents}

\section*{Introduction}
\addcontentsline{toc}{section}{Introduction}

Le chapitre précédent a montré que la Migration Factory se situe à l'intersection de la migration déterministe, de l'assistance intelligente et de l'orchestration multi-agents. Le présent chapitre décrit la conception du système multi-agents (SMA) retenu pour coordonner ces capacités.

La description suit plusieurs points de vue, conformément à ISO/IEC/IEEE 42010: architecture globale, rôles, interactions, états, workflows et gouvernance \cite{iso42010}. Les diagrammes de composants, de séquence, d'activité et d'états utilisent les principes de la notation UML 2.5.1 \cite{omguml}. L'objectif n'est pas de reproduire toutes les classes du code, mais d'exposer les décisions structurantes qui garantissent la modularité, la traçabilité et la maîtrise de l'autorité.

% ============================================================
\section{Choix du modèle SMA}
% ============================================================

\subsection{Architecture hybride et centralement orchestrée}

La Migration Factory adopte une architecture \textbf{multi-agents hybride, spécialisée et centralement orchestrée}. Elle est hybride parce qu'elle associe des agents fondés sur des LLM, des services déterministes, des outils externes et une supervision humaine. Elle est spécialisée parce que chaque agent possède une responsabilité limitée. Elle est centralement orchestrée parce que les transitions, la sélection des tâches et la progression globale sont contrôlées par un orchestrateur et un état persistant.

Ce choix répond à quatre contraintes :

\begin{itemize}
    \item l'ordre des étapes de migration ne peut pas être laissé à une conversation libre entre agents ;
    \item les commandes et écritures doivent rester déterministes ;
    \item les décisions sensibles doivent pouvoir interrompre le workflow ;
    \item les résultats intermédiaires doivent être persistés et vérifiables.
\end{itemize}

\begin{figure}[H]
    \centering
    \resizebox{0.95\textwidth}{!}{%
    \begin{tikzpicture}[
        core/.style={draw=cgired,rounded corners=4pt,fill=white,minimum width=5.3cm,minimum height=1.35cm,align=center,font=\bfseries},
        node/.style={draw=cgiblue,rounded corners=3pt,fill=cgilight!55,text width=4.6cm,minimum height=1.25cm,align=center,font=\small},
        ext/.style={draw=cgigray,rounded corners=3pt,fill=white,text width=4.6cm,minimum height=1.25cm,align=center,font=\small},
        arr/.style={<->,thick,draw=cgigray}
    ]
        \node[core] (orch) {Orchestrateur central\\état, routage, interruptions et reprise};
        \node[node,above left=1.25cm and 1.1cm of orch] (agents) {Agents spécialisés\\analyse, planification, transformation, review et réparation};
        \node[node,above right=1.25cm and 1.1cm of orch] (human) {Supervision humaine\\approbation, rejet, révision et annulation};
        \node[ext,below left=1.25cm and 1.1cm of orch] (services) {Services déterministes\\transitions, commandes, fichiers et validation};
        \node[ext,below right=1.25cm and 1.1cm of orch] (stores) {État et preuves\\base, événements, artefacts et checkpoints};
        \draw[arr] (agents)--(orch);
        \draw[arr] (human)--(orch);
        \draw[arr] (services)--(orch);
        \draw[arr] (stores)--(orch);
    \end{tikzpicture}%
    }
    \caption{Modèle SMA hybride retenu}
    \label{fig:modele-sma-ch3}
\end{figure}

\subsection{Justification des deux dépôts indépendants}

Le SMA est décliné dans deux systèmes indépendants : Java/Spring Boot et Angular. La séparation ne remet pas en cause la cohérence du modèle ; elle permet d'adapter les agents, outils, profils et politiques à chaque écosystème.

\begin{table}[H]
    \centering
    \small
    \arrayrulecolor{cgigray}
    \rowcolors{2}{cgilight!55}{white}
    \begin{tabularx}{\textwidth}{|>{\bfseries\RaggedRight\arraybackslash}p{4.0cm}|>{\RaggedRight\arraybackslash}X|>{\RaggedRight\arraybackslash}X|}
        \hline
        \rowcolor{cgilight!95}
        \textbf{Dimension} & \textbf{Java/Spring Boot} & \textbf{Angular} \\
        \hline
        Runtime & JDK 11, 17 et 21 & Node.js selon la version Angular \\
        \hline
        Gestion des dépendances & Maven et \texttt{pom.xml} & npm, \texttt{package.json} et lockfile \\
        \hline
        Outils de transformation & OpenRewrite et transformations Java/Maven & Angular CLI, \texttt{ng update} et transformations TypeScript \\
        \hline
        Validation & Compilation et tests Maven & Installation, build Angular et tests frontend \\
        \hline
        Rythme d'évolution & Profils Spring Boot/Java & Matrice Angular/Node.js/TypeScript/RxJS \\
        \hline
    \end{tabularx}
    \caption{Spécialisation des deux déclinaisons du SMA}
    \label{tab:specialisation-sma-ch3}
    \rowcolors{2}{white}{white}
    \arrayrulecolor{black}
\end{table}

% ============================================================
\section{Identification et classification des agents}
% ============================================================

Les agents sont regroupés selon la phase à laquelle ils contribuent.

\begin{figure}[H]
    \centering
    \resizebox{0.98\textwidth}{!}{%
    \begin{tikzpicture}[
        group/.style={draw=cgiblue,rounded corners=4pt,fill=cgilight!45,text width=4.2cm,minimum height=3.6cm,align=left,inner sep=8pt,font=\small},
        root/.style={draw=cgired,rounded corners=3pt,fill=white,minimum width=4.5cm,minimum height=1.0cm,align=center,font=\bfseries},
        arr/.style={-{Latex[length=2mm]},thick,draw=cgigray}
    ]
        \node[root] (root) {SMA de migration};
        \node[group,below left=1.2cm and 4.6cm of root] (understand) {\textbf{Comprendre}\\[1mm]
        Agent d'analyse\\
        Reviewer d'analyse\\
        Agent de détection de profil};
        \node[group,below=1.2cm of root] (decide) {\textbf{Décider et organiser}\\[1mm]
        Planificateur\\
        Reviewer du plan\\
        Assistant de décision};
        \node[group,below right=1.2cm and 4.6cm of root] (act) {\textbf{Transformer et vérifier}\\[1mm]
        Agent de transformation\\
        Agent build/test\\
        Proposeur et reviewer de réparation};
        \draw[arr] (root)--(understand);
        \draw[arr] (root)--(decide);
        \draw[arr] (root)--(act);
    \end{tikzpicture}%
    }
    \caption{Classification fonctionnelle des agents}
    \label{fig:classification-agents-ch3}
\end{figure}

\begin{table}[H]
    \centering
    \scriptsize
    \arrayrulecolor{cgigray}
    \rowcolors{2}{cgilight!55}{white}
    \begin{tabularx}{\textwidth}{|>{\bfseries\RaggedRight\arraybackslash}p{2.8cm}|>{\RaggedRight\arraybackslash}p{3.4cm}|>{\RaggedRight\arraybackslash}p{3.5cm}|>{\RaggedRight\arraybackslash}X|}
        \hline
        \rowcolor{cgilight!95}
        \textbf{Agent} & \textbf{Objectif local} & \textbf{Sortie principale} & \textbf{Limite d'autorité} \\
        \hline
        Analyseur & Comprendre le projet et identifier les risques. & Rapport d'analyse et preuves. & Lecture seule sur la source. \\
        \hline
        Reviewer d'analyse & Vérifier la cohérence et la couverture. & Décision structurée et remarques. & Ne modifie ni le projet ni l'état final. \\
        \hline
        Planificateur & Décomposer la cible en unités ordonnées. & Plan de migration. & Le plan doit être revu et approuvé. \\
        \hline
        Reviewer du plan & Contrôler l'ordre, les dépendances et les validations. & Acceptation ou demande de révision. & Ne transforme pas le code. \\
        \hline
        Transformateur & Préparer les modifications attendues. & Changements candidats et artefacts. & Écrit uniquement dans le workspace autorisé. \\
        \hline
        Agent build/test & Déclencher ou interpréter la validation technique. & Résultats, journaux et classification. & Les commandes sont construites par le backend. \\
        \hline
        Proposeur de réparation & Construire un correctif à partir d'un échec. & Diff candidat structuré. & Ne peut ni approuver ni appliquer. \\
        \hline
        Reviewer de réparation & Évaluer et consolider la proposition. & Décision et diff final éventuel. & La sortie reste soumise à l'humain. \\
        \hline
        Assistant & Expliquer l'état et les actions possibles. & Réponse contextualisée. & Aucun pouvoir direct d'exécution. \\
        \hline
    \end{tabularx}
    \caption{Rôles et limites d'autorité des agents}
    \label{tab:roles-agents-ch3}
    \rowcolors{2}{white}{white}
    \arrayrulecolor{black}
\end{table}

% ============================================================
\section{Architecture globale du SMA}
% ============================================================

\subsection{Vue de composants}

\begin{figure}[H]
    \centering
    \resizebox{\textwidth}{!}{%
    \begin{tikzpicture}[
        comp/.style={draw=cgiblue,rounded corners=3pt,fill=cgilight!50,minimum width=3.5cm,minimum height=1.05cm,align=center,font=\small},
        store/.style={draw=cgigray,rounded corners=3pt,fill=white,minimum width=3.5cm,minimum height=1.05cm,align=center,font=\small},
        arr/.style={-{Latex[length=2mm]},thick,draw=cgigray}
    ]
        \node[comp] (ui) {Cockpit et assistant};
        \node[comp,right=0.65cm of ui] (api) {API et services\\applicatifs};
        \node[comp,right=0.65cm of api] (state) {Service d'état et\\transitions};
        \node[comp,right=0.65cm of state] (orch) {Orchestrateur\\LangGraph};

        \node[comp,below=1.0cm of api] (agents) {Agents spécialisés};
        \node[comp,below=1.0cm of state] (executor) {Services déterministes\\et exécuteur};
        \node[store,below=1.0cm of ui] (db) {SQLite\\état autoritaire};
        \node[store,below=1.0cm of orch] (artifacts) {Artefacts, logs\\et workspaces};

        \draw[arr] (ui)--(api);
        \draw[arr] (api)--(state);
        \draw[arr] (state)--(orch);
        \draw[arr] (orch)--(agents);
        \draw[arr] (agents)--(executor);
        \draw[arr] (executor)--(artifacts);
        \draw[arr] (state)--(db);
        \draw[arr,bend right=20] (db) to (api);
        \draw[arr,bend left=20] (artifacts) to (state);
    \end{tikzpicture}%
    }
    \caption{Architecture globale du SMA}
    \label{fig:architecture-sma-ch3}
\end{figure}

La base de données conserve l'état métier autoritaire. LangGraph coordonne le workflow et les interruptions, mais ne remplace pas les services métier. Les agents produisent des artefacts ; les services déterministes vérifient, appliquent et exécutent.

\subsection{Frontières d'autorité}

La conception sépare quatre niveaux :

\begin{enumerate}
    \item \textbf{intention utilisateur :} démarrer, approuver, rejeter, annuler ou demander une révision ;
    \item \textbf{raisonnement agentique :} analyser, planifier, expliquer ou proposer ;
    \item \textbf{décision métier :} vérifier l'état et autoriser une transition ;
    \item \textbf{exécution technique :} construire et lancer une commande ou écrire un fichier.
\end{enumerate}

Un agent ne peut donc pas transformer directement une intention en commande shell. Il produit un résultat intermédiaire, contrôlé par les couches suivantes.

% ============================================================
\section{Coordination et communication}
% ============================================================

\subsection{Modèle de communication}

Les agents ne communiquent pas par une conversation libre persistante. Ils échangent indirectement au moyen d'artefacts et de contrats orchestrés. Chaque message logique contient :

\begin{itemize}
    \item un identifiant de tâche et de migration ;
    \item le rôle du producteur et du consommateur ;
    \item une référence vers le contexte autorisé ;
    \item le schéma de sortie attendu ;
    \item les empreintes des artefacts examinés ;
    \item le statut, les erreurs et les informations de corrélation.
\end{itemize}

Cette organisation reprend le principe des protocoles de communication entre agents, tout en utilisant des contrats applicatifs plutôt que l'ensemble de FIPA-ACL \cite{fipa-acl}.

\subsection{Séquence analyse--planification--transformation}

\begin{figure}[H]
    \centering
    \resizebox{\textwidth}{!}{%
    \begin{tikzpicture}[
        actor/.style={font=\small\bfseries,align=center},
        msg/.style={-{Latex[length=2mm]},thick,draw=cgigray},
        ret/.style={-{Latex[length=2mm]},dashed,draw=cgigray}
    ]
        \node[actor] (dev) at (0,0) {Développeur};
        \node[actor] (orch) at (3.3,0) {Orchestrateur};
        \node[actor] (analyse) at (6.6,0) {Analyseur};
        \node[actor] (review) at (9.9,0) {Reviewer};
        \node[actor] (plan) at (13.2,0) {Planificateur};
        \node[actor] (exec) at (16.5,0) {Exécuteur};
        \foreach \x in {0,3.3,6.6,9.9,13.2,16.5} {\draw[dashed,cgigray] (\x,-0.35)--(\x,-7.1);}
        \draw[msg] (0,-0.9)--node[above,font=\scriptsize]{Lancer}(3.3,-0.9);
        \draw[msg] (3.3,-1.55)--node[above,font=\scriptsize]{Analyser}(6.6,-1.55);
        \draw[ret] (6.6,-2.2)--node[above,font=\scriptsize]{Rapport}(3.3,-2.2);
        \draw[msg] (3.3,-2.85)--node[above,font=\scriptsize]{Réviser}(9.9,-2.85);
        \draw[ret] (9.9,-3.5)--node[above,font=\scriptsize]{Décision}(3.3,-3.5);
        \draw[msg] (3.3,-4.15)--node[above,font=\scriptsize]{Planifier}(13.2,-4.15);
        \draw[ret] (13.2,-4.8)--node[above,font=\scriptsize]{Plan}(3.3,-4.8);
        \draw[ret] (3.3,-5.45)--node[above,font=\scriptsize]{Présenter}(0,-5.45);
        \draw[msg] (0,-6.1)--node[above,font=\scriptsize]{Approuver}(3.3,-6.1);
        \draw[msg] (3.3,-6.75)--node[above,font=\scriptsize]{Transformer et valider}(16.5,-6.75);
    \end{tikzpicture}%
    }
    \caption{Séquence de coordination des agents principaux}
    \label{fig:sequence-agents-ch3}
\end{figure}

% ============================================================
\section{Conception de l'orchestrateur}
% ============================================================

L'orchestrateur reçoit une intention validée par l'API, lit l'état autoritaire puis choisit le prochain nœud. Ses responsabilités sont :

\begin{itemize}
    \item sélectionner l'agent ou le service adapté ;
    \item séquencer les tâches et respecter leurs préconditions ;
    \item transmettre les références d'artefacts ;
    \item enregistrer les transitions et événements ;
    \item interrompre le workflow pour une décision humaine ;
    \item reprendre à partir d'un checkpoint ;
    \item classer les erreurs et choisir entre arrêt, réparation ou replanification ;
    \item empêcher l'exécution d'une transition devenue obsolète.
\end{itemize}

\begin{figure}[H]
    \centering
    \resizebox{0.95\textwidth}{!}{%
    \begin{tikzpicture}[
        step/.style={draw=cgiblue,rounded corners=3pt,fill=cgilight!55,minimum width=3.2cm,minimum height=1.0cm,align=center,font=\small\bfseries},
        decision/.style={draw=cgired,diamond,aspect=2,fill=white,align=center,font=\scriptsize\bfseries},
        arr/.style={-{Latex[length=2mm]},thick,draw=cgigray}
    ]
        \node[step] (receive) {Recevoir l'intention};
        \node[step,right=0.55cm of receive] (load) {Charger l'état};
        \node[decision,right=0.75cm of load] (valid) {Transition\\autorisée ?};
        \node[step,right=0.75cm of valid] (route) {Router la tâche};
        \node[step,right=0.55cm of route] (persist) {Persister le résultat};
        \node[decision,below=1.25cm of route] (next) {Décision ou\\erreur ?};
        \node[step,left=0.75cm of next] (pause) {Interrompre};
        \node[step,right=0.75cm of next] (continue) {Poursuivre};
        \draw[arr] (receive)--(load);
        \draw[arr] (load)--(valid);
        \draw[arr] (valid)--node[above,font=\scriptsize]{oui}(route);
        \draw[arr] (route)--(persist);
        \draw[arr] (persist.south)--(next.north);
        \draw[arr] (next)--node[above,font=\scriptsize]{oui}(pause);
        \draw[arr] (next)--node[above,font=\scriptsize]{non}(continue);
    \end{tikzpicture}%
    }
    \caption{Logique générale de l'orchestrateur}
    \label{fig:orchestrateur-ch3}
\end{figure}

% ============================================================
\section{Gestion des états, de la mémoire et du contexte}
% ============================================================

\subsection{État global et états locaux}

Une migration possède un état global et plusieurs états locaux : stage, commande, point de validation, proposition de réparation et rapport. Le frontend affiche une projection ; il ne détermine pas l'état.

\begin{figure}[H]
    \centering
    \resizebox{0.92\textwidth}{!}{%
    \begin{tikzpicture}[
        state/.style={draw=cgiblue,rounded corners=3pt,fill=cgilight!55,minimum width=2.7cm,minimum height=0.95cm,align=center,font=\small\bfseries},
        terminal/.style={draw=cgired,rounded corners=3pt,fill=white,minimum width=2.7cm,minimum height=0.95cm,align=center,font=\small\bfseries},
        arr/.style={-{Latex[length=2mm]},thick,draw=cgigray}
    ]
        \node[state] (created) {CREATED};
        \node[state,right=0.55cm of created] (ready) {READY};
        \node[state,right=0.55cm of ready] (running) {RUNNING};
        \node[state,right=0.55cm of running] (paused) {PAUSED};
        \node[terminal,right=0.55cm of paused] (complete) {COMPLETED};
        \node[state,below=1.2cm of running] (recovery) {RECOVERY\\REQUIRED};
        \node[terminal,below=1.2cm of paused] (cancelled) {CANCELLED};
        \node[terminal,below=1.2cm of ready] (failed) {FAILED};
        \draw[arr] (created)--(ready);
        \draw[arr] (ready)--(running);
        \draw[arr] (running)--(paused);
        \draw[arr] (paused)--(complete);
        \draw[arr,bend left=25] (paused) to (running);
        \draw[arr] (running)--(recovery);
        \draw[arr] (recovery)--(paused);
        \draw[arr] (paused)--(cancelled);
        \draw[arr] (running)--(failed);
    \end{tikzpicture}%
    }
    \caption{Machine à états simplifiée d'une migration}
    \label{fig:machine-etats-sma-ch3}
\end{figure}

\subsection{Mémoire et contexte}

Trois formes de mémoire sont distinguées :

\begin{itemize}
    \item \textbf{mémoire métier persistante :} états, décisions, événements et références ;
    \item \textbf{mémoire documentaire :} analyses, plans, logs, diffs et rapports ;
    \item \textbf{contexte agentique temporaire :} extraits et informations sélectionnés pour une invocation.
\end{itemize}

Le contexte d'un agent n'est pas l'intégralité du dépôt. Il est construit selon le rôle, borné en taille, nettoyé et lié aux artefacts persistés. Cette stratégie réduit le risque de fuite d'information et améliore la pertinence du raisonnement.

% ============================================================
\section{Workflows SMA spécialisés}
% ============================================================

\subsection{Workflow Java/Spring Boot}

Le workflow Java suit une trajectoire de versions. Pour chaque stage :

\begin{enumerate}
    \item qualifier le JDK, Maven et le profil Spring Boot ;
    \item analyser le \texttt{pom.xml}, le code, la configuration et les tests ;
    \item faire réviser l'analyse ;
    \item construire et faire réviser le plan ;
    \item soumettre le plan au développeur ;
    \item préparer un workspace isolé ;
    \item appliquer les transformations et recettes OpenRewrite ;
    \item exécuter Maven, compiler et tester ;
    \item ouvrir la réparation gouvernée en cas d'échec ;
    \item valider les preuves et passer au stage suivant.
\end{enumerate}

\subsection{Workflow Angular}

Le workflow Angular conserve la même gouvernance, avec des tâches adaptées :

\begin{enumerate}
    \item détecter Angular, Node.js, npm, TypeScript et RxJS ;
    \item résoudre un profil compatible avec la transition majeure ;
    \item créer un workspace sans \texttt{node\_modules}, caches ni sorties générées ;
    \item analyser les manifestes, configurations et sources TypeScript ;
    \item produire puis réviser le plan ;
    \item appliquer les transformations et migrations Angular CLI ;
    \item installer les dépendances ;
    \item exécuter le build et les tests ;
    \item conserver les preuves puis avancer vers la version suivante.
\end{enumerate}

Au moment de la rédaction, les agents principaux Angular sont développés, mais leur intégration complète, la réparation et le reporting restent en consolidation.

% ============================================================
\section{Supervision humaine et points de décision}
% ============================================================

La supervision humaine intervient lorsque le système ne doit pas décider seul : approbation d'un plan, validation d'un correctif, rejet d'une trajectoire ou annulation d'une migration. Une décision est liée à l'état courant et à l'empreinte de l'artefact examiné.

\begin{table}[H]
    \centering
    \small
    \arrayrulecolor{cgigray}
    \rowcolors{2}{cgilight!55}{white}
    \begin{tabularx}{\textwidth}{|>{\bfseries\RaggedRight\arraybackslash}p{3.4cm}|>{\RaggedRight\arraybackslash}p{4.2cm}|>{\RaggedRight\arraybackslash}X|}
        \hline
        \rowcolor{cgilight!95}
        \textbf{Point de décision} & \textbf{Choix possibles} & \textbf{Effet} \\
        \hline
        Analyse ou plan & Approuver, rejeter, demander une révision. & Reprise, replanification ou arrêt. \\
        \hline
        Transformation sensible & Autoriser ou refuser. & Ouverture ou fermeture de l'exécution. \\
        \hline
        Proposition de réparation & Approuver, rejeter, réviser. & Application contrôlée ou nouvelle proposition. \\
        \hline
        Migration active & Annuler. & Arrêt contrôlé et état terminal. \\
        \hline
    \end{tabularx}
    \caption{Principaux points de décision humaine}
    \label{tab:decisions-humaines-ch3}
    \rowcolors{2}{white}{white}
    \arrayrulecolor{black}
\end{table}

% ============================================================
\section{Boucle de réparation gouvernée}
% ============================================================

\begin{figure}[H]
    \centering
    \resizebox{0.98\textwidth}{!}{%
    \begin{tikzpicture}[
        step/.style={draw=cgiblue,rounded corners=3pt,fill=cgilight!55,minimum width=3.0cm,minimum height=1.0cm,align=center,font=\small\bfseries},
        human/.style={draw=cgired,rounded corners=3pt,fill=white,minimum width=3.0cm,minimum height=1.0cm,align=center,font=\small\bfseries},
        arr/.style={-{Latex[length=2mm]},thick,draw=cgigray}
    ]
        \node[step] (failure) {Conserver l'échec};
        \node[step,right=0.4cm of failure] (context) {Construire le contexte};
        \node[step,right=0.4cm of context] (propose) {Proposer un diff};
        \node[step,right=0.4cm of propose] (review) {Réviser le diff};
        \node[human,right=0.4cm of review] (approve) {Décision humaine};
        \node[step,below=1.15cm of approve] (apply) {Appliquer};
        \node[step,left=0.4cm of apply] (test) {Recompiler et tester};
        \node[step,left=0.4cm of test] (result) {Valider ou réessayer};
        \draw[arr] (failure)--(context);
        \draw[arr] (context)--(propose);
        \draw[arr] (propose)--(review);
        \draw[arr] (review)--(approve);
        \draw[arr] (approve)--(apply);
        \draw[arr] (apply)--(test);
        \draw[arr] (test)--(result);
        \draw[arr,bend left=28] (result.north) to (context.south);
        \draw[arr,bend left=25] (approve.north) to (propose.north);
    \end{tikzpicture}%
    }
    \caption{Boucle de réparation gouvernée du SMA}
    \label{fig:reparation-sma-ch3}
\end{figure}

Les invariants sont les suivants : contexte borné, schéma structuré, diff contrôlé, approbation humaine, application dans le workspace, revalidation obligatoire et nombre limité de tentatives. Le reviewer peut consolider la proposition, mais ne peut pas l'appliquer.

% ============================================================
\section{Gestion des erreurs et replanification}
% ============================================================

Les erreurs sont classées en quatre catégories :

\begin{itemize}
    \item \textbf{précondition invalide :} runtime, chemin ou profil absent ;
    \item \textbf{erreur technique transitoire :} timeout, interruption ou indisponibilité ;
    \item \textbf{échec de transformation ou de validation :} compilation, build ou test ;
    \item \textbf{incohérence fonctionnelle :} analyse ou plan incomplet.
\end{itemize}

Selon la catégorie, l'orchestrateur bloque, réessaie, ouvre la réparation, demande une décision ou renvoie vers la planification. La replanification ne modifie pas silencieusement un plan approuvé : elle crée une nouvelle version soumise à review et décision.

\section*{Conclusion}
\addcontentsline{toc}{section}{Conclusion}

Ce chapitre a défini une architecture SMA adaptée aux contraintes de la migration logicielle. Le système repose sur des agents spécialisés, une orchestration centrale, un état persistant, des services déterministes et une supervision humaine. Les interactions s'effectuent par contrats et artefacts plutôt que par une conversation libre, ce qui permet de contrôler les responsabilités et de relier chaque décision à une preuve.

La séparation des dépôts Java/Spring Boot et Angular conserve les spécificités techniques sans abandonner le modèle commun. Le chapitre suivant décrit l'implémentation concrète de cette architecture : technologies, structure des projets, développement des agents, orchestration, backend, persistance, interface, sécurité et difficultés techniques.
